\documentclass[11pt]{article}
\usepackage[utf8]{inputenc}
\usepackage{geometry}
\geometry{letterpaper, margin=1in}
\usepackage{xcolor}
\usepackage{enumitem}
\usepackage{hyperref}
\usepackage[scaled]{helvet}
\renewcommand{\familydefault}{\sfdefault}

\title{\textbf{Response to Reviewers}}
\author{Manuscript ID: JCM-20140}
\date{\today}

\begin{document}

\maketitle

\noindent \textbf{Manuscript Title:} Enhancing QoS in Dense IEEE 802.11ax Networks using a Dynamic Airtime-Based Soft Admission Control Mechanism \\
\textbf{Authors:} Dayanand Ambawade, Rohan Pawar \\
\textbf{Journal:} Journal of Communications

\vspace{0.5cm}

\noindent Dear Editor and Reviewers,

\vspace{0.3cm}

We sincerely thank the editorial team and the reviewers for their insightful and constructive feedback on our manuscript. We have carefully considered all comments and have revised the paper accordingly.

\noindent As requested, all significant changes in the revised manuscript are highlighted in \textcolor{blue}{blue text}. Below is a point-by-point response to each reviewer's comments.

\vspace{0.5cm}
\hrule
\vspace{0.5cm}

\section*{Response to Reviewer 1}

\noindent \textbf{Comment 1.1:} \textit{Network Health Metrics: Currently, only PER is used as an example. It is recommended to discuss the applicability of other potential metrics (e.g., channel busy/idle ratio, retransmission rate), or to elaborate on the rationale for choosing PER and its practicality/accessibility in real-world devices.}

\noindent \textbf{Response:} Thank you for this excellent suggestion. We have added a new paragraph in \textbf{Section IV.C} explaining our rationale for prioritizing Packet Error Rate (PER) over other metrics like CCA (Clear Channel Assessment). We clarify that PER is an ``end-to-end'' metric that captures hidden node collisions and actual goodput degradation, whereas CCA might underestimate congestion in dense deployments with efficient spatial reuse.

\vspace{0.3cm}

\noindent \textbf{Comment 1.2:} \textit{Overhead and Scalability: The statement "only requires over-the-air timing estimation" is mentioned but not quantified... A brief analysis of the computational complexity of AS-CAC+ on the AP and its potential impact on system performance is suggested.}

\noindent \textbf{Response:} We agree that a quantifiable complexity analysis strengthens the paper. We have added a new subsection \textbf{VII.B (Complexity and Scalability Analysis)} in the Discussion. We explicitly state that the AS-CAC+ decision logic operates in $O(1)$ time per admission request, as it relies on simple arithmetic comparisons, making it significantly lighter than $O(N \cdot M)$ machine learning inference operations required by competing DRL approaches.

\vspace{0.3cm}

\noindent \textbf{Comment 1.3:} \textit{Comparison with Related Work... It could be further enhanced by briefly discussing under what network scales or levels of dynamics/churn the simpler AS-CAC+ might be superior or inferior to more complex learning-based approaches.}

\noindent \textbf{Response:} We have expanded the \textbf{Comparison with Related Work} in Section VII to address this. We now explicitly mention that AS-CAC+ offers superior responsiveness in ``high-churn'' environments (rapid user arrival/departure) where ML models typically suffer from convergence delays or retraining lag.

\vspace{0.5cm}
\hrule
\vspace{0.5cm}

\section*{Response to Reviewer 2}

\noindent \textbf{Comment 2.1:} \textit{Clarity and conciseness: Some sections, particularly the Discussion and Results analysis, could be slightly condensed to avoid repetition of similar observations across figures and tables.}

\noindent \textbf{Response:} We accepted this feedback and have significantly condensed the \textbf{Discussion (Section VII.A)} segment. Instead of repeating numerical findings from the Results section, we have merged the ``Key Findings'' and ``Throughput Gain'' analysis into a single, concise evaluation summary labeled ``Condensed Performance Analysis.''

\vspace{0.3cm}

\noindent \textbf{Comment 2.2:} \textit{Threshold sensitivity: While threshold values are well motivated, a brief sensitivity analysis or discussion on how robust these values are across different environments would strengthen the generality of the approach.}

\noindent \textbf{Response:} We have added a new subsection \textbf{VII.C (Parameter Sensitivity and Robustness)}. We explain that the adaptive nature of AS-CAC+ makes the system resilient to initial parameter choices; even if the starting thresholds are suboptimal, the feedback loop automatically corrects them within seconds.

\vspace{0.3cm}

\noindent \textbf{Comment 2.3:} \textit{Real-world validation: The ns-3 evaluation is thorough, but a short discussion on expected challenges in real hardware or firmware implementations would further improve practicality.}

\noindent \textbf{Response:} We have added subsection \textbf{VII.D (Real-World Implementation Challenges)}, identifying key practical hurdles such as driver-level access to microsecond-granularity airtime counters and the need to synchronize admission decisions with TWT schedules.

\vspace{0.3cm}

\noindent \textbf{Comment 2.4:} \textit{Minor presentation issues: A small number of grammatical and formatting inconsistencies are present and can be addressed with light editing.}

\noindent \textbf{Response:} We have conducted a thorough proofread of the entire manuscript to correct grammatical errors and improve flow and consistency.

\vspace{0.5cm}

\noindent We believe that these revisions have significantly strengthened the quality and clarity of the manuscript. We look forward to your positive response.

\vspace{1cm}

\noindent Sincerely,\\
\textbf{Dayanand Ambawade, Rohan Pawar}\\
Sardar Patel Institute of Technology (SPIT)\\
Mumbai, India

\end{document}
